% META: {"id": "TEXP-COMPLEXES-TRIGO-POLYNOMES-RE-C1", "chapitre": "TEXP-COMPLEXES-TRIGO-POLYNOMES", "type_objet": "remediation", "capacites_codes": ["C1"], "status": "generated"}

%---------------------------------------------------------------
% FICHE DE REMEDIATION — C1 : passer d'une forme a l'autre
%---------------------------------------------------------------

\subsection*{Rappel -- Formes algebrique, trigonometrique, exponentielle}

\textbf{Erreur 1 :} calculer l'argument par $\arctan\left(\frac{b}{a}\right)$
sans regarder le quadrant. Pour $z=-1+\mathrm{i}$, cette formule donne
$-\frac{\pi}{4}$ alors que l'argument est $\frac{3\pi}{4}$ : le point est dans
le deuxième quadrant.

\textbf{Erreur 2 :} oublier le module en repassant à la forme algébrique.
$3\mathrm{e}^{\mathrm{i}\pi/6}$ vaut
$3\cos\frac{\pi}{6}+3\mathrm{i}\sin\frac{\pi}{6}$, pas
$\cos\frac{\pi}{6}+\mathrm{i}\sin\frac{\pi}{6}$.

\textbf{Erreur 3 :} additionner les modules dans un produit. Dans un produit
les modules se \emph{multiplient} et seuls les arguments s'additionnent.

% BEGIN-VERIFY
% from sympy import I, sqrt, Abs, arg, pi, exp, simplify, Rational
% z = -1 + I
% assert arg(z) == 3 * pi / 4
% assert arg(z) != -pi / 4
% w = 3 * exp(I * pi / 6)
% assert simplify(w.rewrite('cos') - (3 * sqrt(3) / 2 + Rational(3, 2) * I)) == 0
% assert simplify(Abs(z * w) - 3 * sqrt(2)) == 0
% assert Abs(z * w) != Abs(z) + Abs(w)
% END-VERIFY

\begin{exercice}{TEXP-COMPLEXES-TRIGO-POLYNOMES-RE-C1-EX1}{1}{10}
\begin{enumerate}
  \item Déterminer un argument de $z=-1+\mathrm{i}$ en plaçant le point, puis
        expliquer pourquoi $\arctan(-1)$ ne donne pas la bonne valeur.
  \item Calculer $|z\times 3\mathrm{e}^{\mathrm{i}\pi/6}|$ et vérifier que ce
        n'est pas $|z|+3$.
\end{enumerate}
\end{exercice}
